\chapter{Synthetic extension spectral sequences}
\label{chap:synthetic-extensions}

This chapter covers Section 4.  We first transport the general extension
spectral sequence weight by weight, then show that the connecting map in a
$\lambda$-quotient triangle encodes precisely the classical Adams
differentials.  The finite-quotient statement below makes explicit an image
that is ambiguous in the paper's printed Proposition 4.6.

\section{Weightwise extension spectral sequences}

% SOURCE: aimpaper/main.tex:998-1007.
\begin{definition}[Synthetic extension spectral sequence]
  \label{def:synthetic-extension-ss}
  \uses{def:synthetic-adams-ss,def:f-extension-ss}
  \notready
  For a weight-preserving map $f:X\to Y$ of finite synthetic spectra with
  strongly convergent synthetic Adams spectral sequences, filter
  $0\to\pi_{*,*}X\to\pi_{*,*}Y\to0$ by Adams filtration.  The resulting
  $f$-ESS has
  \[
    {}^fE_0^{s,t,w}\cong E_\infty^{s,t,w}(X)\oplus
    E_\infty^{s,t,w}(Y),
    \qquad d_n^f:(s,t,w)\mapsto(s+n,t+n,w).
  \]
\end{definition}

% SOURCE: aimpaper/main.tex:999-1007, assertion that all Section 2 results remain valid.
\begin{proposition}[Weightwise transport of the ESS calculus]
  \label{prop:synthetic-ess-weightwise-transport}
  \uses{def:synthetic-extension-ss,prop:fextension-iff-detected,prop:fextension-inessential-iff-higher,prop:no-crossing-iff-uniform-detection,thm:fess-commutative-square,prop:fess-exact-middle-boundary}
  \notready
  Every definition and theorem of Chapter~\ref{chap:extension-ss} holds in a
  fixed synthetic weight, provided every map and comparison preserves that
  weight.  In particular, essentiality, crossings, the representative
  criterion, square propagation, and exact-middle boundary detection retain
  their statements with a third grading coordinate.
\end{proposition}

\begin{proof}
  Fix $w$ and regard the bigraded homotopy groups in that weight as filtered
  graded groups.  Reuse the two-term filtered-complex construction and each
  proof from Chapter~\ref{chap:extension-ss}; all maps preserve $w$, so no
  additional reindexing occurs.  Reassemble the family over all weights.
\end{proof}

% SOURCE: aimpaper/main.tex:1009-1021, Notation not:deltaandrho.
\begin{definition}[$\lambda$--$\rho$--$\delta$ quotient triangle]
  \label{def:lambda-rho-delta-triangle}
  \uses{def:synthetic-lambda-quotient,def:cofiber-triangle}
  \notready
  For $0<n<m\le\infty$, define the distinguished triangle
  \[
    \Sigma^{0,-n}\nu X/\lambda^{m-n}
      \xrightarrow{\lambda^n}\nu X/\lambda^m
      \xrightarrow{\rho_{n,m}}\nu X/\lambda^n
      \xrightarrow{\delta_{n,m}}
      \Sigma^{1,-n}\nu X/\lambda^{m-n}.
  \]
  For $m=\infty$, the first two terms are
  $\Sigma^{0,-n}\nu X$ and $\nu X$.  Translation in the synthetic category is
  suspension by $S^{1,0}$.
\end{definition}

% SOURCE: BHS Proposition A.13; aimpaper/main.tex:1016-1019.
\begin{theorem}[$\lambda$-adic completeness]
  \label{thm:external-synthetic-lambda-complete}
  \uses{def:lambda-rho-delta-triangle,def:external-result}
  \notready
  An explicit BHS external result supplies a natural equivalence
  \[
    \nu X\simeq\varprojlim_m\nu X/\lambda^m
  \]
  compatible with the quotient maps and the finite
  $\lambda$--$\rho$--$\delta$ triangles.
\end{theorem}

\begin{proof}
  Instantiate BHS Proposition A.13 and transport its grading and inverse-system
  maps to the conventions of Definition~\ref{def:lambda-rho-delta-triangle}.
\end{proof}

% SOURCE: aimpaper/main.tex:1981-2081, representative equations used in the
% page-restriction argument; project packaging of their inverse system.
\begin{definition}[Finite ESS solution torsors and coherent towers]
  \label{def:synthetic-ess-solution-tower}
  \uses{def:synthetic-extension-ss,def:fess-cycles-boundaries,cor:fess-map-of-stable-pages}
  \notready
  Let $F_m:A_m\to B_m$, $m\ge m_0$, be compatible maps of finite
  $\lambda$-quotients, with restriction maps from level $m+1$ to level $m$.
  Fix a synthetic weight, an extension length $n$, and compatible source and
  target classes $x_m,y_m$.  The \emph{solution fiber}
  $\operatorname{Sol}_m^n(x_m,y_m)$ consists of representative-level
  witnesses for
  \[
    d_n^{F_m}(x_m)=y_m:
  \]
  a filtered representative of $x_m$, a representative of the prescribed
  complete target coset of $y_m$, and a sum of shorter $F_m$-ESS images whose
  representative equation realizes that coset.  Two witnesses are identified
  exactly when their source representatives differ by a filtered boundary and
  their target equations differ by the same shorter-image correction.

  When nonempty, $\operatorname{Sol}_m^n(x_m,y_m)$ is a torsor under the
  abelian \emph{difference group} $K_m$ of changes preserving the source class
  and complete target coset.  Stable-page naturality gives equivariant maps
  \[
    \operatorname{res}_{m+1,m}:\operatorname{Sol}_{m+1}^n
      \longrightarrow\operatorname{Sol}_m^n,
    \qquad K_{m+1}\longrightarrow K_m.
  \]
  A \emph{coherent solution tower} is a sequence
  $u_m\in\operatorname{Sol}_m^n$ satisfying
  $\operatorname{res}_{m+1,m}(u_{m+1})=u_m$ for every $m\ge m_0$.
\end{definition}

\begin{proof}
  Use the filtered two-term complex defining the synthetic ESS.  A choice of
  filtered source lift and target correction is precisely a solution of its
  representative equation.  Subtracting two solutions gives the stated
  difference group, and addition acts freely and transitively.  The compatible
  quotient squares induce the restriction maps and their equivariance.
\end{proof}

% PROOF-SOURCE: the standard obstruction calculation for a countable inverse
% system of abelian torsors, combined with BHS Proposition A.13.
\begin{lemma}[Mittag--Leffler criterion for an inverse-limit ESS extension]
  \label{lem:synthetic-ess-coherent-tower-limit}
  \uses{def:synthetic-ess-solution-tower,thm:external-synthetic-lambda-complete}
  \notready
  In the setting of
  Definition~\ref{def:synthetic-ess-solution-tower}, suppose every finite
  solution fiber is nonempty and that the compatible maps $F_m$ are the
  reductions of a map $F_\infty$ between $\lambda$-complete synthetic spectra.
  There is a canonical torsor obstruction
  \[
    \omega(\operatorname{Sol}_\bullet)
      \in \varprojlim^{1}_{m} K_m
  \]
  whose vanishing is equivalent to the existence of a coherent solution
  tower.  Either of the following is sufficient for vanishing:
  \begin{enumerate}
    \item every finite solution lifts, that is, each
      $\operatorname{res}_{m+1,m}$ is surjective;
    \item the difference-group tower $(K_m)$ is Mittag--Leffler: for every
      $m$, the images of $K_l\to K_m$ stabilize as $l$ increases.  Surjective
      transition maps are a sufficient special case.  Then
      $\varprojlim^1_m K_m=0$.
  \end{enumerate}
  A coherent solution tower determines the untruncated relation
  \[
    d_n^{F_\infty}(x_\infty)=y_\infty
  \]
  in the complete target coset.  Conversely, every untruncated representative
  solution restricts to a coherent finite tower.
\end{lemma}

\begin{proof}
  Choose one point in each finite solution torsor.  The differences between a
  chosen point and the restriction of the next one form the usual cocycle for
  the inverse system $(K_m)$; changing the choices changes it by a coboundary.
  Its class is $\omega$, and a coboundary correction makes the chosen points
  coherent exactly when $\omega=0$.  Surjective lifting constructs such a
  tower recursively.  The Mittag--Leffler hypothesis gives
  $\varprojlim^1K_m=0$ and hence the same conclusion.

  At the representative-space level, the solution equation is a homotopy fiber
  of the filtered two-term-complex map, and homotopy limits preserve that
  fiber.  The $\lambda$-adic equivalences identify the limit source, target,
  map, and shorter-image corrections with their untruncated counterparts.
  Thus a coherent finite solution is exactly an $F_\infty$ solution; restriction
  proves the converse.
\end{proof}

\section{The elementary $\lambda$ and $\rho$ extensions}

% SOURCE: aimpaper/main.tex:1025-1048, Proposition prop:9770ae6e.
\begin{proposition}[$\lambda^n$- and $\rho$-ESS have only $d_0$]
  \label{prop:lambda-rho-ess-only-d0}
  \uses{def:lambda-rho-delta-triangle,prop:synthetic-einfty-map-compatibility,prop:synthetic-einfty-first-quotient,prop:synthetic-ess-weightwise-transport,prop:fess-exact-middle-boundary}
  \notready
  In the extension spectral sequences of $\lambda^n$ and $\rho_{n,m}$, the
  only possibly nonzero differentials are
  $d_0^{\lambda^n}=\lambda^n$ and $d_0^\rho=\rho$.  These $d_0$ extensions have
  no crossings.  On each tridegree, $\lambda^n$ is surjective or zero and
  $\rho$ is injective or zero on $E_\infty$.
\end{proposition}

\begin{proof}
  First suppose $m<\infty$.  If one of the quotient exponents is one, use the
  mod-$\lambda$ collapse and unique-detection statement of
  Proposition~\ref{prop:synthetic-einfty-first-quotient}; when both relevant
  exponents are at least two, use the two finite $E_\infty$ formulas.  If a
  tridegree is outside the quotient range its group is zero, while in range
  $\rho$ is a cycle inclusion and $\lambda^n$ is a boundary quotient.

  If $m=\infty$, instead use the two infinite-endpoint clauses of
  Proposition~\ref{prop:synthetic-einfty-map-compatibility}: on the
  untruncated $E_\infty$ formula, $\lambda^n$ is the quotient that raises the
  boundary cutoff by $n$, and
  $\rho:\nu X\to\nu X/\lambda^n$ includes permanent representatives into the
  finite cycle cutoff (using the special-fiber clause when $n=1$).  Thus the
  sequence is exact at its middle term in every tridegree in both the finite
  and infinite branches.
  Apply the synthetic exact-middle theorem to eliminate higher differentials,
  and use the degree-zero filtration shift to obtain no-crossing.
\end{proof}

\section{The connecting extension and classical differentials}

% SOURCE: aimpaper/main.tex:1068-1131, proof of Proposition prop:6de7d130.
\begin{lemma}[Shorter $\delta_n$ images are classical boundaries]
  \label{lem:delta-images-equal-classical-boundaries}
  \uses{thm:external-lambda-bockstein,prop:synthetic-einfty-map-compatibility,prop:synthetic-ess-weightwise-transport,cor:fess-triangle-shift}
  \notready
  For the infinite connecting map
  $\delta_n:\nu X/\lambda^n\to\Sigma^{1,-n}\nu X$, the subgroup generated by
  extension differentials of length less than $r$ in the target tridegree is
  exactly the image of the classical boundary subgroup
  $B_{r-1}^{s+r,t+r-1}(X)$ under the corresponding $\lambda$-multiple.
\end{lemma}

\begin{proof}
  Induct on $r$.  Write each generator of $B_{r-1}$ as the target of an
  essential classical $d_{r'}$ with $2\le r'<r$.  The Bockstein comparison and
  triangle-shift propagation turn it into a shorter $\delta_n$ extension with
  the required $\lambda$ power.  Conversely, apply the same comparison to
  every shorter $\delta_n$ extension.  The inductive hypothesis removes its
  lower-page indeterminacy.
\end{proof}

% SOURCE: aimpaper/main.tex:1068-1131, Proposition prop:6de7d130, m=infinity case.
\begin{proposition}[Infinite $\delta_n$-ESS formula]
  \label{prop:delta-ess-formula-infinite}
  \uses{lem:delta-images-equal-classical-boundaries,thm:external-lambda-bockstein,cor:fess-triangle-shift,prop:lambda-rho-ess-only-d0}
  \notready
  Suppose the classical Adams spectral sequence has
  $d_r(x)=y$, with $x\in Z_{r-1}^{s,t}$ and
  $y\in Z_\infty^{s+r,t+r-1}$.  For
  $\delta_n:\nu X/\lambda^n\to\Sigma^{1,-n}\nu X$:
  \begin{enumerate}
    \item if $r\ge n+1$, then
      $d_r^{\delta_n}(x)=\lambda^{r-n-1}y$;
    \item if $r\le n+1$, then
      $d_r^{\delta_n}(\lambda^{n+1-r}x)=y$.
  \end{enumerate}
  The formulas agree when $r=n+1$ and are equalities in the correct ESS
  quotient, whether or not the displayed extension is essential.
\end{proposition}

\begin{proof}
  Start with $n=1$, where the assertion is the $\lambda$-Bockstein theorem.
  Induct on $n$ using the commuting triangle relating multiplication by
  $\lambda$, $\delta_{n-1}$, and $\delta_n$.  Triangle propagation gives the
  formula before quotient indeterminacy.  Division by one $\lambda$ introduces
  only the next classical boundary subgroup; Lemma~\ref{lem:delta-images-equal-classical-boundaries}
  identifies and removes precisely that ESS indeterminacy.
\end{proof}

% SOURCE: aimpaper/main.tex:1068-1131, Proposition prop:6de7d130, finite m case.
% VERIFIED-IN: the printed second case places y in E_infty(nu X); the finite target requires its rho-image.
\begin{proposition}[Finite $\delta_{n,m}$-ESS formula]
  \label{prop:delta-ess-formula-finite}
  \uses{prop:delta-ess-formula-infinite,cor:fess-composition,prop:lambda-rho-ess-only-d0}
  \notready
  For finite $m>n$, compose the formulas of
  Proposition~\ref{prop:delta-ess-formula-infinite} with the reduction
  $\rho:\nu X\to\nu X/\lambda^{m-n}$.  Thus the displayed target is the image
  of $\lambda^{r-n-1}y$, or of $y$ in the second case, in the finite quotient;
  it is zero once the required $\lambda$ power lies outside that quotient.
  This explicit image corrects the type ambiguity in
  \texttt{aimpaper/main.tex:1088--1094}.
\end{proposition}

\begin{proof}
  Factor $\delta_{n,m}$ as $\delta_{n,\infty}$ followed by $\rho$.  The
  $\rho$-ESS has only its injective-or-zero $d_0$, with no crossing, so ESS
  composition transports the infinite formula to the finite target.  The
  quotient formula decides exactly when the image is zero.
\end{proof}

% SOURCE: aimpaper/main.tex:1138-1144, Corollary cor:2a636737.
\begin{corollary}[$\lambda$-multiple form of the $\delta$ formula]
  \label{cor:delta-ess-lambda-multiples}
  \uses{prop:delta-ess-formula-finite}
  \notready
  In the situation above,
  \[
    d_r^\delta(\lambda^a x)=\lambda^{a+r-n-1}y
  \]
  whenever $0\le a\le n$ and
  $0\le a+r-n-1<m-n$; outside the finite quotient range the right-hand side
  and the extension are zero.
\end{corollary}

\begin{proof}
  Multiply the finite formula by $\lambda^a$, using $\lambda$-linearity of the
  synthetic Adams and extension spectral sequences.  Check the source and
  target quotient ranges directly.
\end{proof}

% SOURCE: aimpaper/main.tex:1146-1150, Remark rmk:qvoewfj.
\begin{theorem}[$\delta$-ESS and classical differential cosets are equivalent]
  \label{thm:delta-classical-equivalence}
  \uses{cor:delta-ess-lambda-multiples,lem:delta-images-equal-classical-boundaries,thm:external-lambda-bockstein}
  \notready
  For every allowed $a,n,m,r$, the equation
  $d_r^\delta(\lambda^a x)=\lambda^{a+r-n-1}y$ holds if and only if the
  classical differential coset is $d_r(x)=y$.  Both sides have the same
  indeterminacy $B_{r-1}$ after applying the stated $\lambda$ power.
\end{theorem}

\begin{proof}
  The forward direction applies the Bockstein comparison to a representative
  and uses the shorter-image lemma to identify the quotient.  The reverse
  direction is the $\lambda$-multiple formula.  Equality of the two boundary
  subgroups makes the correspondence bijective at the level of differential
  cosets, not merely at the level of nonvanishing.
\end{proof}

% SOURCE: aimpaper/main.tex:1152-1163, Example 4.13.
\begin{proposition}[$\delta_1,\delta_2,\delta_3$ regression]
  \label{prop:delta-14-stem-regression}
  \uses{prop:synthetic-14-stem-regression,cor:delta-ess-lambda-multiples}
  \notready
  The $d_2(h_4)=h_0h_3^2$ and
  $d_3(h_0h_4)=h_0d_0$ examples yield exactly the eight
  $\delta_1,\delta_2,\delta_3$ extensions and $\lambda$ powers listed in
  Example 4.13: two for $\delta_1$, three for $\delta_2$, and three for
  $\delta_3$.
\end{proposition}

\begin{proof}
  Substitute $(n,m)=(1,\infty),(2,\infty),(3,\infty)$ and each admissible
  value of $a$ into Corollary~\ref{cor:delta-ess-lambda-multiples}.  The
  synthetic $E_\infty$ regression verifies that every source and target exists
  in the displayed tridegree.
\end{proof}

\section{Crossings of classical differentials}

% SOURCE: aimpaper/main.tex:1165-1173, Definition def:classicalcrossdiff.
\begin{definition}[Classical crossing on an $E_{n+1}$ page]
  \label{def:classical-differential-crossing}
  \uses{def:classical-adams-ss,def:ss-cycle-boundary-permanent}
  \notready
  Let $r\ge n+1$ and $d_r(x)=y$, with $x$ in filtration $s$.  A crossing on
  the $E_{n+1}$ page is an essential differential
  $d_{r-a-b}(x')=y'$ where the source filtration is $s+a$, the target
  filtration is $s+r-b$, $0<a\le n-1$, and
  $0\le b\le r-n-1$.
\end{definition}

% SOURCE: aimpaper/main.tex:1218-1234, Proposition prop:cross-dr-En.
\begin{proposition}[Classical crossing iff $\delta_n$ crossing]
  \label{prop:classical-crossing-iff-delta-crossing}
  \uses{def:classical-differential-crossing,def:fess-crossing,thm:delta-classical-equivalence,thm:external-synthetic-einfty-nu,thm:external-synthetic-einfty-quotient,prop:synthetic-einfty-first-quotient}
  \notready
  A classical $d_r(x)=y$ has a crossing on $E_{n+1}$ if and only if the
  corresponding extension
  $d_r^{\delta_n}(x)=\lambda^{r-n-1}y$ has a crossing in the synthetic
  $\delta_n$-ESS.
\end{proposition}

\begin{proof}
  A synthetic crossing has the unique degree form
  $d_{r-a-b}^{\delta_n}(\lambda^a x')=\lambda^{r-n-1-b}y'$.  The two
  sides are both empty when $n=1$: the classical range would require
  $0<a\le0$, while mod-$\lambda$ collapse and unique detection exclude an
  additional synthetic representative.  For $n\ge2$, the two $E_\infty$
  formulas force exactly the inequalities in
  Definition~\ref{def:classical-differential-crossing}.  Apply the
  $\delta$--classical equivalence in each direction, preserving essentiality
  through the shared boundary quotient.  The first-quotient proposition is
  used only in the $n=1$ branch.
\end{proof}

% SOURCE: aimpaper/main.tex:1236-1238, Remark nocrossE2.
\begin{corollary}[No classical crossing on $E_2$]
  \label{cor:classical-e2-no-crossing}
  \uses{def:classical-differential-crossing,prop:classical-crossing-iff-delta-crossing}
  \notready
  No classical Adams differential has a crossing on the $E_2$ page, and every
  corresponding $\delta_1$ extension has no crossing.
\end{corollary}

\begin{proof}
  Setting $n=1$ would require an integer $a$ with $0<a\le0$.  Hence no
  classical crossing exists; transfer the conclusion through
  Proposition~\ref{prop:classical-crossing-iff-delta-crossing}.
\end{proof}

% SOURCE: aimpaper/main.tex:1240-1271, Example exam:classcrossdiff.
\begin{proposition}[Stem-$38$ crossing regression]
  \label{prop:stem38-crossing-regression}
  \uses{prop:classical-crossing-iff-delta-crossing,cor:classical-e2-no-crossing,def:external-evidence}
  \notready
  Given the typed computation evidence
  $d_3(e_1)=h_1t$ and
  $d_4(h_0h_3h_5)=h_0^2x$, the former is a crossing of the latter on the
  $E_3$ and $E_4$ pages, but not on $E_2$ or on pages $E_r$ for $r\ge5$.
  The associated $\delta_n$ extensions and their weights are exactly those in
  Example 4.18.
\end{proposition}

\begin{proof}
  Check the four inequalities in the crossing definition for $n=1,2,3$ and
  $n\ge4$, then transport the two computed differential cosets using the
  crossing equivalence.  The synthetic $E_\infty$ formulas verify which
  $\lambda$ multiples remain represented.
\end{proof}
